% SPAR Midterm Report Template
% Compatible with arXiv (TeX Live 2023+, pdflatex)
%
% To compile: latexmk -pdf report.tex
% Or use the build scripts: build report.tex

% ============================================================================
% DOCUMENT CLASS
% ============================================================================
% Options: midterm | final | arxiv
% - midterm: For mid-program report
% - final: For end-of-program report
% - arxiv: For converting your report into a preprint for arXiv
\documentclass[midterm]{sparreport}

% ============================================================================
% METADATA
% ============================================================================
\title{Paper Gains, Realized Losses: Testing the Realization Effect in Large Language Models}

\author{
  \textbf{Ciaran Walsh}\\
  \texttt{crw2161@columbia.edu} \\
  Columbia University
  \and
  \textbf{Emilio Barkett\footnote{SPAR Mentor}}\\
  \texttt{eab2291@columbia.edu}\\
  Columbia University
}
\sparround{Fall 2025}

% ============================================================================
% DOCUMENT BEGINS
% ============================================================================
\begin{document}

\maketitle

% ============================================================================
% ABSTRACT
% ============================================================================
\begin{abstract}
This project tests whether large language models (LLMs) reproduce the realization
effect documented in human casino gamblers by Flepp, Meier, and Franck \cite{flepp-2021-effect}.
We adapt their casino framework into a vignette experiment in which prior outcomes
are manipulated to be either paper (within-visit, mental account open) or realized
(between-visits, mental account closed), and measure subsequent wager choices.
Across nine models and three sampling temperatures, we estimate condition effects
using OLS regression with model, temperature, and prompt-version fixed effects and
HC3 robust standard errors. Preliminary pooled results show substantial and
statistically significant condition effects, but the signs of many coefficients
diverge from the human-field predictions: paper losses are associated with
\emph{reduced} wagering relative to baseline, and realized losses with
\emph{increased} wagering, reversing the realization effect. The one partial
exception is the magnitude ordering within paper gains, where large gains produce
higher wagers than small gains, consistent with H2b. These reversals motivate
further analysis at the model level and across alternative prompt framings.
\end{abstract}

% Keywords
\keywords{realization effect \and large language models \and behavioral economics \and risk-taking \and mental accounting}

% ============================================================================
% INTRODUCTION
% ============================================================================
\section{Introduction}

The realization effect, formalized by Imas (2016) and empirically validated by
Flepp et al.\ \cite{flepp-2021-effect} predicts that people take more risk after paper outcomes than
after realized outcomes. When a prior loss or gain remains unrealized---that is,
when the gambler has not yet cashed out and closed their mental account---it stays
integrated with future prospects, motivating loss-chasing behavior or encouraging
further risk-taking on the cushion of paper gains. Once outcomes are realized
through a money transfer, the mental account closes, the reference point resets,
and risk behavior returns toward baseline or, for large losses, falls below it.

This question is significant for AI behavioral modeling. Large language models are
increasingly used in decision-support tools, preference elicitation, risk
simulation, and policy prototyping. If LLM behavior systematically diverges from
human decision patterns in well-documented ways, downstream applications that
assume model-human behavioral correspondence may produce misleading outputs.
Conversely, if LLMs do replicate human biases, this raises its own concerns: models
trained on human-generated text may inherit and amplify behavioral heuristics in
contexts where rational calculation is warranted.

We replicate the Flepp et al.\ \cite{flepp-2021-effect} field study design as a vignette experiment,
presenting LLM agents with slot machine scenarios that vary both prior
outcome magnitude and account status (paper versus realized). We test the four
directional hypotheses from the original paper---that paper losses and gains each
increase risk-taking, that larger outcomes amplify these effects, that realized
losses decrease risk-taking, and that realized gains leave risk-taking unchanged---
and examine whether LLM behavior matches, mirrors, or is orthogonal to the
human-field pattern.

% ============================================================================
% METHODOLOGY
% ============================================================================
\section{Methodology}

\subsection{Experimental Design}

We implement an 11-condition vignette experiment structured around the quintile
design of Flepp et al. (2021, Table 2). Each condition specifies a prior outcome
amount in Swiss francs (CHF) and an account status. Paper conditions present the
agent as currently inside the casino with the balance remaining on a playing card;
realized conditions present the agent as returning for a new visit after having
cashed out at the end of a prior visit. This framing operationalizes the open versus
closed mental account distinction from Imas (2016).

Paper outcome conditions use \texttt{paper\_even} (CHF 0) as the baseline,
corresponding to the paper's Q3 quintile (--96 to 0 CHF). Non-baseline paper
conditions cover three loss levels---small (--60 CHF, within Q3), medium (--200
CHF, Q2: --309 to --97), and large (--350 CHF, Q1: $\leq$ --310)---and two gain
levels---small (+40 CHF, Q4: 1--80) and large (+150 CHF, Q5: $\geq$ 81).

Realized outcome conditions use \texttt{realized\_small\_loss} (--30 CHF) as the
baseline (Q4: --62 to 0 CHF). Non-baseline realized conditions cover three loss
levels---medium (--400 CHF, Q3: --787 to --63), large (--1800 CHF, Q2: --2790 to
--788), and extreme (--3500 CHF, Q1: $\leq$ --2791)---and one gain level (+200
CHF, Q5: $\geq$ 1). The medium realized loss and realized gain conditions were not
statistically significant in the original paper, providing useful null-effect
benchmarks.

\subsection{Stimuli}

Each trial presents a system prompt establishing the casino visitor role and a
condition-specific user prompt. The primary prompt version (\texttt{absolute})
states prior outcomes as explicit CHF amounts. A second version (\texttt{balance})
frames outcomes as card balance relative to the start of the visit. A third version
(\texttt{qualitative}) describes outcomes in relative terms only (e.g., ``a modest
amount'', ``a substantial amount'') with no numeric figures, testing whether the
realization effect holds when the reference point is implicit rather than
numerically specified. The key manipulation is consistent across all versions:
paper scenarios specify the balance remains on the playing card, while realized
scenarios specify the player has cashed out and returned for a new visit.

Each trial elicits two responses: (1) a total session wager in CHF (range 1--1000,
reflecting the paper's mean session wager of approximately 650 CHF), and (2) a
slot machine risk preference on a 1--5 scale (1 = low-volatility, 5 = high-volatility),
providing an analog to the paper's second dependent variable (\texttt{LogT-winCasino}).

\subsection{Models and Parameters}

We query models via the OpenRouter API. Current results include nine models spanning
several major providers and capability tiers: \texttt{anthropic/claude-3.5-haiku},
\texttt{anthropic/claude-haiku-4.5}, \texttt{deepseek/deepseek-chat-v3-0324},
\texttt{google/gemini-3.1-flash-lite-preview}, \texttt{gpt-4.1-mini},
\texttt{openai/gpt-5.4-mini}, \texttt{qwen/qwen3-32b},
\texttt{moonshotai/kimi-k2-thinking}, and \texttt{x-ai/grok-4-fast}. Each
condition is run at three sampling temperatures (0.5, 1.0, 1.5) to assess
sensitivity to output stochasticity.

\subsection{Analysis}

We estimate OLS regressions of log-wager on condition dummies, with model,
temperature, and prompt-version included as covariate fixed effects and HC3
heteroscedasticity-robust standard errors throughout. Paper and realized conditions
are estimated in separate regressions, as each trial tests one outcome type only.
Hypothesis tests mirror H1a--H4 from Flepp et al., with one-sided $p$-values
where the original paper specifies a directional prediction, and Wald tests on
coefficient differences for the magnitude-ordering hypotheses (H1b, H2b, H3b).

% ============================================================================
% PRELIMINARY RESULTS
% ============================================================================
\section{Preliminary Results}

Current data comprise 12,155 trials, of which 12,138 yield valid parsed wagers,
across nine models and three temperatures. We first estimate OLS regressions of
log-wager on outcome conditions.

\begin{table}[h]
\centering
\begin{tabular}{lcc}
\hline
 & Coefficient & (SE) \\
\hline
\multicolumn{3}{c}{\textit{Panel A: Paper Outcomes (Baseline: Paper Even)}} \\
\hline
Paper Loss (Small)  & -0.1833*** & (0.0224) \\
Paper Loss (Medium) & -0.5085*** & (0.0177) \\
Paper Loss (Large)  & -0.3584*** & (0.0192) \\
Paper Gain (Small)  & -0.5024*** & (0.0166) \\
Paper Gain (Large)  &  0.1231*** & (0.0165) \\
\hline
Observations & \multicolumn{2}{c}{6,065} \\
$R^2$ & \multicolumn{2}{c}{0.495} \\
\hline
\\[-1em]
\multicolumn{3}{c}{\textit{Panel B: Realized Outcomes (Baseline: Small Realized Loss)}} \\
\hline
Realized Loss (Medium)  & 0.3948*** & (0.0245) \\
Realized Loss (Large)   & 0.4653*** & (0.0232) \\
Realized Loss (Extreme) & 0.4175*** & (0.0236) \\
Realized Gain           & 0.0786*** & (0.0158) \\
\hline
Observations & \multicolumn{2}{c}{6,073} \\
$R^2$ & \multicolumn{2}{c}{0.448} \\
\hline
\end{tabular}
\caption{OLS estimates of log-wager on outcome conditions. Regressions include
fixed effects for model, temperature, and prompt framing. HC3 robust standard
errors in parentheses. *** $p < 0.01$.}
\end{table}

\begin{table}[h]
\centering
\begin{tabular}{lcc}
\hline
Hypothesis & Prediction & Result \\
\hline
H1a & Paper losses $\uparrow$ risk & Not supported \\
H1b & Larger losses $\uparrow$ risk more & Partially supported \\
H2a & Paper gains $\uparrow$ risk & Not supported \\
H2b & Large gain $>$ small gain & Supported \\
H3a & Realized losses $\downarrow$ risk & Not supported \\
H3b & Larger losses $\downarrow$ risk more & Not supported \\
H4  & Realized gain $=$ baseline & Rejected \\
\hline
\end{tabular}
\caption{Summary of hypothesis tests relative to Flepp et al. (2021). Results
are based on pooled preliminary estimates and should be interpreted with caution
given the data quality issues described in Section 4.}
\end{table}

Preliminary pooled estimates show that LLM behavior does not straightforwardly
reproduce the human realization-effect pattern and in several respects reverses it.

In Panel A, all paper loss coefficients are negative relative to the paper-even
baseline, indicating that LLMs reduce their wager following paper losses rather
than increase it as predicted. The ordering is also non-monotone: the medium loss
(--0.51) produces a larger reduction than the large loss (--0.36), partially
violating H1b. Paper gains are similarly discordant: small paper gains reduce
wagering (--0.50), while large paper gains increase it (+0.12), making H2b the
only hypothesis to receive clear support.

In Panel B, all realized loss coefficients are positive---LLMs increase wagering
after a realized loss relative to the small-loss baseline, reversing H3a. The
magnitude ordering is approximately monotone for the medium-to-large comparison
(+0.39 vs. +0.47), though the extreme loss slightly reverses (+0.42), providing
only weak support for H3b. Realized gains also significantly increase wagering
(+0.08), rejecting H4.

One interpretation of these reversals is that models are responding to the
numerical magnitude of stated outcomes---treating a large stated loss as a signal
to be cautious---rather than to the paper-versus-realized framing. The asymmetric
pattern for gains (small reduces, large increases) is consistent with a
magnitude-based reference point: small gains may anchor the model to a modest
outcome and suppress wagering, while large gains license greater risk-taking
regardless of account status. The near-identical coefficients for \texttt{paper\_loss\_small}
and \texttt{paper\_loss\_medium} in the preliminary data also reflect a data quality
issue described below.

% ============================================================================
% CHALLENGES & OPEN QUESTIONS
% ============================================================================
\section{Challenges \& Open Questions}

\paragraph{Between-conditions vs.\ within-person identification.}
The most fundamental structural difference from Flepp et al. (2021) is that the
original study uses individual-level fixed-effects OLS, comparing the \emph{same
gambler's} behavior across different prior-outcome states. This controls for all
stable individual differences in risk tolerance. Our vignette design is purely
between-conditions: each LLM trial is an independent query with no individual
gambling history to exploit. The closest analog---model and temperature fixed
effects---absorbs systematic differences across models but cannot replicate the
within-person identification of the field study. This is unavoidable in a vignette
design and situates our study closer to the Imas (2016) laboratory replication
than the Flepp et al. field study.

\paragraph{Legacy condition labels in existing data.}
The current results file contains data collected across multiple experimental
iterations with inconsistent condition naming. Early runs used \texttt{paper\_neutral}
and \texttt{realized\_neutral} for the baseline conditions (now
\texttt{paper\_even} and \texttt{realized\_small\_loss}), and
\texttt{realized\_loss\_medium}, \texttt{realized\_loss\_large}, and
\texttt{realized\_loss\_small} for what are now \texttt{realized\_medium\_loss},
\texttt{realized\_large\_loss}, and \texttt{realized\_small\_loss}. Critically,
\texttt{paper\_loss\_small} in the existing data contains runs at --200 CHF (the
value previously associated with the mislabeled condition, now renamed
\texttt{paper\_loss\_medium}), whereas the corrected \texttt{paper\_loss\_small}
is defined at --60 CHF. Pooled analysis that does not account for this mixing will
produce biased condition estimates. The preliminary results above should be read
with this caveat.

\paragraph{Wager range and prompt framing.}
Initial data collection used a 1--500 CHF response range and prompt wording that
was ambiguous between a per-spin and a session-total wager interpretation. The
500 CHF cap binds for a non-trivial fraction of trials given the paper's mean
session wager of approximately 650 CHF, artificially truncating the upper tail of
the response distribution and attenuating detectable loss-chasing effects. Updated
prompts raise the cap to 1--1000 CHF and clarify that responses represent total
session wagers.

\paragraph{Explicit numeric reference points.}
The \texttt{absolute} prompt version presents exact CHF figures to the model. Real
casino customers do not necessarily track their running balance with this
precision---they experience losses and gains but may not consciously anchor on a
specific number. Providing an exact amount may cause LLMs to apply more deliberate,
arithmetic reasoning rather than simulating affective loss-chasing, which could
explain the conservative responses to paper losses. The \texttt{qualitative} prompt
version, which uses only relative descriptors, will help isolate whether numeric
precision is driving the divergence from human-field results.

\paragraph{Pooling across heterogeneous models.}
The current dataset is heavily imbalanced: \texttt{gpt-4.1-mini} contributes
approximately 7,000 of the 12,155 rows, representing 58\% of the sample. Pooled
coefficients are therefore dominated by one model's behavioral patterns. Whether
the reversals documented above are a property of all models or specific to
particular architectures or capability tiers is an open question that
model-level analysis will address.

% ============================================================================
% NEXT STEPS
% ============================================================================
\section{Next Steps}

\paragraph{Data cleaning and harmonization.}
The immediate priority is to reconcile legacy condition labels in the existing
results file so that the --200 CHF \texttt{paper\_loss\_small} rows are correctly
treated as \texttt{paper\_loss\_medium}, and baseline conditions are consistently
identified regardless of naming variant. A cleaned, consolidated dataset will allow
model-level and prompt-version-level subgroup analyses without confounding from
labeling inconsistencies.

\paragraph{Per-model analysis.}
Once data are cleaned, we will rerun the regression analysis separately for each of
the nine models using \texttt{analyze\_results.py --per-model}. This will establish
whether the pooled reversals are uniform across models or driven by a subset,
whether reasoning models (e.g., \texttt{kimi-k2-thinking}) diverge systematically
from non-reasoning models, and whether effects scale with model capability tier.

\paragraph{New data collection under the updated design.}
Fresh runs using the updated prompts (1--1000 CHF range, explicit session-total
framing, two-response format capturing both wager and slot machine risk preference)
will replace or supplement the truncated legacy data. We will also collect data
under the \texttt{qualitative} prompt version, which has not yet been run, to test
whether removing numeric anchors brings LLM behavior closer to the human-field
pattern.

\paragraph{Expanded model coverage.}
We will extend collection to additional models from our evaluation list, prioritizing
frontier and reasoning models not yet covered. Of particular interest are models
that have been studied in prior behavioral economics replications
(e.g., \texttt{openai/gpt-4o}, \texttt{openai/o3}) to enable comparison with
the emerging literature on LLM decision-making.

\paragraph{Risk-profile dependent variable.}
Updated trials elicit a slot machine risk preference (1--5) alongside the wager
amount, providing a second dependent variable analogous to the paper's
\texttt{LogT-winCasino}. If wager size and machine choice covary across conditions
in a coherent direction, this will strengthen the interpretation of whatever
realization-effect pattern emerges; if they dissociate, it will suggest that
different response channels are governed by different mechanisms within the models.

\paragraph{Prompt-version comparison.}
The final analysis will compare results across \texttt{absolute},
\texttt{balance}, and \texttt{qualitative} prompt versions. If the realization
effect is present under \texttt{qualitative} framing but absent or reversed under
\texttt{absolute} framing, this would indicate that LLMs apply explicit numerical
reasoning that overrides the mental-accounting heuristics the paper-versus-realized
manipulation is designed to elicit. This would be a substantive finding in its own
right regarding the conditions under which LLM behavior mirrors human behavioral
patterns.

% ============================================================================
% BIBLIOGRAPHY
% ============================================================================
\bibliographystyle{plain}
\bibliography{references}

\end{document}
